\documentclass[conference]{IEEEtran}
\usepackage{amsmath,amssymb}
\usepackage{booktabs}
\usepackage[hidelinks]{hyperref}

\begin{document}

\title{Kernel Concentration and Bandwidth Scaling\\in Quantum Kernel Methods: Project Synopsis}

\author{\IEEEauthorblockN{Muhammad Ahmad Masood and Muhammad Waqas}
\IEEEauthorblockA{AI 631 Quantum Machine Learning, Fall 2026 --- Project \#5}}

\maketitle

\begin{abstract}
Fidelity quantum kernels can concentrate: as the number of qubits grows, off-diagonal kernel values cluster around a fixed value, and estimating them with a finite number of measurement shots can no longer separate different inputs. Scaling the input features by a bandwidth parameter has been proposed as a way to counter this. We will measure how fast the off-diagonal kernel variance decays from 2 to 16 qubits as a function of bandwidth. We will also determine when that variance falls below the shot-noise floor, and test whether a tuned bandwidth keeps classification accuracy competitive with an equally tuned classical RBF kernel. Experiments use exact statevector simulation, finite-shot sampling, a noise model derived from IBM device calibration data, and runs on IBM quantum hardware at two to three qubit counts.
\end{abstract}

\section{Research Question and Hypotheses}
\textbf{Research question.} How fast do quantum kernel values concentrate as the qubit count grows, and does a bandwidth hyperparameter restore learnability?

For a data-encoding unitary $U(\cdot)$ and bandwidth $\lambda>0$, we study the fidelity kernel
\begin{equation}
k_\lambda(\mathbf{x},\mathbf{x}') = \bigl|\langle 0^{n}|\,U^{\dagger}(\lambda\mathbf{x}')\,U(\lambda\mathbf{x})\,|0^{n}\rangle\bigr|^{2},
\end{equation}
where $\mathbf{x}\in\mathbb{R}^{n}$ uses one feature per qubit. Our central quantity is the off-diagonal variance $\sigma^{2}_{\lambda}(n)=\mathrm{Var}_{i\neq j}\,[k_\lambda(\mathbf{x}_i,\mathbf{x}_j)]$ over the data. When $k$ is estimated as the frequency of the all-zero outcome in $N$ shots of a compute--uncompute circuit, each entry carries binomial noise with variance $k(1-k)/N$.

\textbf{H1 (concentration).} At fixed $\lambda=1$, $\sigma^{2}_{\lambda}(n)$ decreases exponentially in $n$, as predicted for expressive encodings~\cite{thanasilp2024concentration}.

\textbf{H2 (bandwidth).} There is a decreasing bandwidth schedule $\lambda^{*}(n)$ under which $\sigma^{2}_{\lambda^{*}}(n)$ stays above the shot-noise level at a fixed shot budget, and test accuracy at $n=16$ stays statistically above that obtained at $\lambda=1$, consistent with~\cite{shaydulin2022bandwidth,canatar2022bandwidth}.

\textbf{Falsifier.} We reject H2 in our setting if, at $n=16$ and $N=10^{4}$ shots, every $\lambda$ in the search grid gives an off-diagonal spread that cannot be distinguished from the shot-noise spread (variance-ratio test, $\alpha=0.05$), \emph{or} if the best tuned $\lambda$ does not give significantly higher test accuracy than $\lambda=1$ across fixed outer folds. Small $\lambda$ is not automatically a fix: as $\lambda\to0$ every kernel value tends to $1$, which is another form of concentration. The test therefore looks for an intermediate regime.

\section{Background}
Quantum kernel methods embed data into quantum states and use state overlaps as a kernel for a classical support vector machine~\cite{havlicek2019supervised,schuld2019feature}. Huang \emph{et al.}\ showed that when quantum kernel values for distinct inputs become very small, the kernel matrix approaches the identity and generalization suffers~\cite{huang2021power}. Thanasilp \emph{et al.}\ identified expressivity, entanglement, global measurements, and hardware noise as sources of exponential concentration. Under concentration, a polynomial number of shots cannot resolve kernel values, and the resulting models become insensitive to the input data~\cite{thanasilp2024concentration}. Shaydulin and Wild showed numerically that tuning a feature-scaling bandwidth can restore generalization for quantum kernels with many qubits~\cite{shaydulin2022bandwidth}, and Canatar \emph{et al.}\ analyzed how bandwidth controls generalization~\cite{canatar2022bandwidth}.

\section{Data and Models}
\textbf{Datasets.} (i) Fashion-MNIST~\cite{xiao2017fashion}, restricted to one binary class pair (T-shirt/top vs.\ Pullover), class-balanced and subsampled to a few hundred images per run, with five subsampling seeds. (ii) A controlled synthetic dataset, a two-class Gaussian mixture in $\mathbb{R}^{n}$ with tunable separation, which isolates the effect of $n$ and $\lambda$ from image structure.

\textbf{Preprocessing.} Standardization and PCA to $n$ components are fit on the training indices of each fold only, followed by rescaling the features to a fixed angle range before applying $\lambda$. The course NB11 starter notebook was not available to us. We therefore implement equivalent fixed folds, leakage-free preprocessing, baselines, ablations, and noise sweeps ourselves, and we test explicitly that no transform sees test data.

\textbf{Quantum model.} An IQP-type $ZZ$ feature map~\cite{havlicek2019supervised} with nearest-neighbour entanglement and a small number of repetitions, applied to $\lambda\mathbf{x}$. The precomputed kernel is passed to an SVM. We sweep $n\in\{2,4,\dots,16\}$ and $\lambda$ over a logarithmic grid. For every circuit we will report the logical and transpiled depth and the two-qubit gate count.

\textbf{Simulation regimes.} (a) Exact statevector kernels. (b) Finite-shot kernels obtained by binomial sampling with $N\in\{10^{2},\dots,10^{5}\}$. (c) Noisy simulation with a noise model built from an IBM backend's calibration data, including a sweep that scales the error rates. (d) Hardware.

\textbf{Hardware.} Compute--uncompute circuits for a small set of data points at two to three qubit counts (tentatively $n\in\{2,4,6\}$) and two bandwidths. The circuits are transpiled for, and run on, an IBM quantum processor. The same transpiled circuits, data points, and shot counts are then simulated with a noise model built from that backend's calibration at run time, so that the off-diagonal variances can be compared directly. We will record the backend name, job IDs, shots, transpilation settings, and calibration data. Our IBM Quantum account (Open plan) provides 600~s of QPU time per 28-day period, so the number of circuits and shots will be sized to fit within that allocation.

\section{Baselines and Evaluation}
All models share fixed, stratified outer folds (5 folds) stored on disk, with inner cross-validation for model selection. Baselines: an RBF-kernel SVM whose $\gamma$ is searched on a logarithmic grid of the same size and with the same selection protocol as $\lambda$, a linear SVM, logistic regression, and the majority-class rate. We report the mean and standard deviation of test accuracy over outer folds and seeds, off-diagonal variance with bootstrap confidence intervals, and fitted decay rates of $\sigma^{2}_{\lambda}(n)$. We will report cases where the classical baseline outperforms the quantum kernel.

\section{Execution Plan}
\begin{table}[h]
\centering
\caption{Week-by-week plan}
\begin{tabular}{@{}p{0.23\columnwidth}p{0.70\columnwidth}@{}}
\toprule
Due & Deliverables \\
\midrule
Synopsis & This document; software environment; IBM account access (done). \\
Nov 15 & Data pipeline, fixed folds, leakage tests, classical baselines, first statevector kernel; Introduction and Methods drafts; code on GitHub. \\
Nov 22 & Full 2--16 qubit sweep; finite-shot analysis; noise-model runs and noise sweep; ablations (feature-map repetitions, entangling vs.\ product map, training-set size); draft Results. \\
Nov 29 & IBM hardware runs and matched noise-model comparison; slides. \\
Dec 6 & Final report (PDF and \LaTeX{} source), tagged GitHub release, presentation. \\
\bottomrule
\end{tabular}
\end{table}

\textbf{Main risks.} Queue time and limited QPU allocation for the hardware runs, which we address by submitting early and keeping the circuit count small. Simulation cost at $n=16$, which we address with statevector inner products rather than sampled circuits for the exact kernels.

\bibliographystyle{IEEEtran}
\bibliography{refs}

\end{document}
